\chapter{Literature review: information markets and scalable oversight}\label{ch:litrev-oversight}
% Assembled from the related-work sections of the papers in Part II
% (litrev/infonomy.tex, litrev/solib.tex). Lifted text is verbatim;
% composed glue is draft-marked.

\begin{draftblock}
\noindent The chapters of this part draw on two literatures: (1) the economics of information, which studies how information can be valued and traded, and (2) the AI-alignment literature on scalable oversight, which studies how to get correct reward signals for strong agents from weak overseers. We first review work in the first area, then turn to the second, and finally review the empirical work on scalable oversight protocols that \cref{ch:solib} builds on.
\end{draftblock}

\section{The economics of information}\label{litrev:sec:infoecon}
\label{infonomy:sec:related}
\paragraph{Information economics.} The foundational theory of information economics\bleg{feel like this might elicit some annoying comments from reviewers} is the \emph{value-of information} framework \citeps{kuhn1953information, howard1966information, raiffa1961AppliedStatisticalDecisionTheory}, which treats information as an instrumental good \citeps{vanalstyneProposalValuingInformation1999}. The buyer's inspection paradox was introduced by \citets{arrow1972economic} and named by \citets{vanalstyneProposalValuingInformation1999}. \citets{hansonIPBarbedWire2011,hansonRahEfficientIP2011} discussed the inefficiency of information markets in the context of intellectual property (IP) law, commenting: ``\emph{just as farmers developed barbed-wire, someday I expect IP advocates will develop better forms of intellectual property}''.

\paragraph{Mechanism design for information markets.} Simple, naive information markets suffer from a number of flaws: \emph{the low cost of duplication} \citeps{samuelson2009economics}, \emph{the transaction costs of tenders} \citeps{tenders}, the transaction cost of \emph{learning} new information, and \emph{information asymmetry} \citeps{arrow1972economic,vanalstyneProposalValuingInformation1999}. \citets{conitzerPredictionMarketsMechanism2012} introduced a mechanism for rewarding information-providing agents based on their influence on prediction market prices, though this is only applicable in contexts where ground-truth (prediction market resolutions) is available.
\paragraph{Zero-knowledge proofs.} Another way for a seller to prove the value of his information without revealing it is via a \emph{zero-knowledge proof} \citeps{goldreich2001FoundationsCryptographyV1} -- in formal settings, this is available if the information is a solution to a PSPACE problem \citeps{zkPSPACE1,zkPSPACE2}; extending this to informal settings is an active area of work \citeps{hammondNeuralInteractiveProofs2024}.

\paragraph{Miscellaneous} Recent works like \citets{manelli2006Bundling,babiaoff2012OptimalMechanismsSellingInformation,bergemannDesignPriceInformation2018} have studied ``optimal mechanisms for selling information'', but from the point-of-view of a seller maximizing his revenue -- we, on the other hand, are interested in improving the efficiency of the information market itself. Information market mechanisms have also been designed for \emph{data markets} in machine learning, e.g. \citets{fallah2024ThreeLayerDataMarkets,chenSellingDataMachine2022,ghorbani2019DataShapley}.\bleg{all papers in this para are totally irrelevant but I suspect reviewers may bring them up}

\section{Mechanism design with language-model agents}\label{litrev:sec:llm-mechanisms}
\paragraph{Mechanism design with LLMs.} Market design for LLM participants has opened up many new frontiers previously not possible with humans alone. Apart from the information bazaar \citeps{weissRedesigningInformationMarkets2024}, this includes e.g. token auctions for online advertising \citeps{duttingMechanismDesignLarge2024}, economic simulations with AI agents \citeps{liEconAgentLargeLanguage2024a}, and forecasting with LLMs \citeps{schoeneggerWisdomSiliconCrowd2024,halawiApproachingHumanLevelForecasting2024,palekaConsistencyChecksLanguage2024,buterin2024InfoFinance}.

\section{Scalable oversight}\label{litrev:sec:oversight}
\label{solib:sec:related}
\paragraph{Scalable oversight} The fundamental limitation of RLHF that it relies on a human's ability to judge a (potentially superhuman) AI's outputs has long been recognized \citeps{christiano2018IDA}, and is known as the \emph{scalable oversight} problem in the AI alignment literature \citeps{hubinger2020alignmentProposalsComplexityClasses,bowmanMeasuringProgressScalable2022}. One well-known scalable oversight proposal is Debate \citeps{irvingAISafetyDebate2018}; our proposal to augment RLHF with information markets may be seen as yet another such proposal.
\paragraph{Oversight protocols.} Apart from Debate, proposed protocols for scalable oversight include: \emph{Iterated Amplification} \citep{christiano2018IDA}, \emph{market-making} \citep{evanhubingerAISafetyMarket2020}, \emph{self-critique} \citep{saunders2022selfcritique}, \emph{reward-modelling} \citep{leike2018RewardModeling} and proposed improvements to Debate such as \emph{doubly-efficient debate} \citep{brown-cohenScalableAISafety2023}. A slightly dated review and discussion of these can be found in \citet{bowmanMeasuringProgressScalable2022}. 

\paragraph{Weak-to-strong generalization and human feedback.} Scalable oversight can be seen as an approach to \emph{weak-to-strong generalization} \citep{sangWeak2Strong, langWeak2Strong} that explicitly relies on the weak model (or human) providing reward to a strong model (as opposed to e.g. fine-tuning or transfer learning). The relationship between scalable oversight and human feedback is made explicit by e.g. \citet{cheng2024RLHFDebate}, who consider \emph{reinforcement learning from debate feedback}.

\paragraph{Previous Debate Experiments.} The following works: \allpreviousworks, all apply an empirical lens to the scalable oversight problem (specifically Debate), similar to our work. While there are important differences in the experiments they conduct, their methodology can broadly be described as measuring the effect of Debate on \emph{judge accuracy}, relative to a baseline of ``Consultancy'' (a protocol where the agent AI is randomly assigned an answer to argue for, and does not have an adversary). % and subtleties

\draft{The recursive information markets of \cref{ch:infonomy} belong to the same family as market making \citep{evanhubingerAISafetyMarket2020}, which they generalize beyond binary questions, and the benchmark of \cref{ch:solib} replaces the judge-accuracy methodology of the debate experiments above with a measure of the incentive a protocol gives for truthfulness.}
